\begin{defn}
  A \tdef{dynamical system} is a tuple $(\tset ,\base ,p\colon\tset\to\base ,f\colon\tset\to\tset )$, where $\tset$ and $\base$ are topological spaces, and $p$ and $f$ are continuous maps such that
  \[
    \begin{tikzcd}
      \tset\arrow[dr,swap,"p"] \arrow[rr,"f"] & & \tset\arrow[dl,"p"]\\
      &\base
    \end{tikzcd}
  \]
  commutes. We also call the $(X,p,f)$ a dynamical system over $\base$. 
\end{defn}
We think of $f$ as the evolution of the dynamical system and $\base$ as a space of parameters. For example, if $\base = \RR$ and $p_1\colon\RR\times\tset\to\RR$ is the projection onto the first coordinate, then a dynamical system \[(\RR\times \tset ,\,\RR ,\, p_1,\, f )\] as above is uniquely determined by the ``flow'' map \[\phi\colon\RR\times\tset\too\tset,\quad (t,x)\mapsto p_2f(t,x),\]
where $p_2\colon\RR\times\tset\to\tset$ denotes the projection onto the second coordinate. (Note however that there is in general no reason for $\phi$ to be compatible with the group structure on $\RR$.)

We will consider dynamical systems satisfying a mild technical condition:

\begin{defn}
    We call a dynamical system $(\tset ,p, f)$ over $\base$ 
    \tdef{shapely} if the map $p\colon X\to \base$ is a proper shape submersion.
\end{defn}


\subsection{Parametrized traces and the Nielsen fixed point theorem}

Let us begin with the following construction:

\begin{cons}
    Let $\psTop{\base}$ be the full subcategory of 
    $\pTop{\base}$ spanned by those $X\to \base$ which are shape
    submersions.
    For any $X\in \psTop{\base}$ and any (small enough) $U\subseteq \base$, the topological space $\Map_{/\base}(U,X)$ has constant
    shape. Thus, we may define the \tdef{relative shape} 
    of $X$ as the sheaf of spaces $\mdef{\ani{X}[\base]}\in \pSpc{\base}$
    given by the sheafification of the presheaf of spaces
    $\ani{\Map_{\Top_{/\base}}(-,X)} \in \Fun(\mathrm{Open}(\base)^\op,\Spc)$. 
    \todo{Find citations and make sure this is actually right}
    The relative shape construction defines a functor
    \[ \ani{-}[\base] \colon \psTop{\base} \too \pSpc{\base}\,.\]
\end{cons}

Consequently, given any shapely dynamical system $(X,p,f)$  
over $\base$, we have an attached sheaf of spaces $\ani{X}[\base]$
together with an endomorphism which we also denote by 
$f\colon \ani{X}[\base]\to \ani{X}[\base]$.

In particular, for a shapely dynamical system 
$f\curvearrowright \tset\xrightarrow p\base$, we may apply 
\cref{cons: Lefschetz index} as well as \cref{cons: Reidemeister trace}
to define the associated Lefschetz index and Reidemeister trace
\begin{equation}
  \label{eq:nielsenthy-intro-characters}
  \lambda^f\colon \cptSec{X^f}[\base]\too \SS_\base\,
  \qquad\text{and}\qquad\rho_{f}\colon \SS_\base \too \SS [\cL^{f}\ani{\tset}[\base]]\,.
\end{equation}
The goal of this subsection is to establish the following relationship between $\lambda^f$ and $\rho_{f}$:
\begin{thm}\label{thm: lefschetz-nielsen}
  Let $f\curvearrowright \tset\xrightarrow p\base$ be a shapely dynamical system. There is a commutative diagram in the category of sheaves of spectra on $\base$:
  \[\begin{tikzcd}
      &  \pbmsec{X^f}{\base}\\
      \SS_\base & {\SS [\cL^{f}\ani{\tset}[\base]]},
      \arrow["\mathfrak a", from=1-2, to=2-2]
      \arrow["{\lambda_f}", from=2-1, to=1-2]
      \arrow["{\rho_f}"', from=2-1, to=2-2]
    \end{tikzcd}\]
  where $\lambda_f$ denotes the $\SS_\base$-linear dual of $\lambda^f$. Furthermore,
  \begin{enumerate}[label=(\roman*)]
  \item $\mathfrak a$ is natural in proper maps of shapely dynamical systems; and
  \item for $f = \id$, the map $\mathfrak a$ is equivalent to the composition
    \[
      \pbmsec{X^f}{\base}\xtoo\sim\SS\lbrack \ani{\tset}[\base]\rbrack\xtoo{\hphantom{\sim}} \SS\lbrack \cL \ani{\tset}[\base]\rbrack,
    \]
    where the first map is the canonical equivalence and the second map is induced by the fiberwise ``constant loop'' map $\ani{\tset}[\base]\to \cL \ani{\tset}[\base]$, which is adjoint to the fiberwise projection $S_\base^1\times \ani{\tset}[\base]\to \ani{\tset}[\base]$.
  \end{enumerate}
\end{thm}
As an immediate consequence of this result, we will deduce the classical Nielsen fixed point theorem, see~\cref{subsec:classical-nielsen}.

The proof of~\cref{thm: lefschetz-nielsen} follows the same pattern as our treatment of the Lefschetz trace formula in the preceding section. Since $\rho_f$ can be defined as the trace of the lax commutative square
\begin{equation}
  \label{eq:reidemeister lax square}
  \begin{tikzcd}
    \Sp_\base\arrow[d,equal] \arrow[r] & \Sp_\base^{\ani{\tset}[\base]}\arrow[d,"f^*"]\\
    \Sp_\base\arrow[r]\arrow[ru,Rightarrow, shorten <=5pt,shorten >=5pt] & \Sp_\base^{\ani{\tset}[\base]}
  \end{tikzcd},
\end{equation}
our goal is to factor this lax commutative square through a traceable endomorphism that encodes the topology of the dynamical system. Such a factorization is provided by the assembly functors and Wall objects of Bartels--Nikolaus and Volpe--Wolff--Ramzi, as we explain in the next section.
\subsubsection{Factoring the Reidemeister trace}\label{subsec:factoring-reidemeister}
A key ingredient in our proof of~\cref{thm: lefschetz-nielsen} is the following bit of general nonsense:
\begin{lem}[Bartels--Nikolaus]\label{lem: bartels-nikolaus nonsense}
  Let $\cC $ be a symmetric monoidal category and let $A$, $B$, and $C$ be objects in $\cC$. Suppose we are given the data of a copairing and a pairing of the form
  \[
    \delta\colon\mathbf 1\too A\otimes B\quad\text{and}\quad
    \nu\colon B\otimes C\too\mathbf 1.
  \]
  Let
  \[
    \mdef{a_{\delta ,\nu}}\colon C\too A
  \]
  denote the composition
  \[
    C\simeq \mathbf 1\otimes C\xtoo{\delta\otimes\id} A\otimes B\otimes C\xtoo{\id\otimes\nu}A\otimes\mathbf 1\simeq A.
  \]
  Given a symmetric monoidal functor $F\colon\cC\to\cD$ where $\cD$ is closed, there is a canonical factorization
  \[\begin{tikzcd}
      &  F(B)^\vee\\
      F(C) & {F(A)},
      \arrow["{a_{F(\delta ),\varepsilon }}", from=1-2, to=2-2]
      \arrow[from=2-1, to=1-2]
      \arrow["{F(a_{\delta ,\nu})}"', from=2-1, to=2-2]
    \end{tikzcd}\]
  natural in $\delta$ and $\nu$,
  where $\varepsilon\colon F(B)^\vee\otimes F(B)\to\mathbf 1$ is the evaluation pairing in $\cD$ and $F(C)\to F(B)^\vee$ is adjoint to the pairing $F(C)\otimes F(B)\simeq F(C\otimes B)\rar{F(\nu )} F(\mathbf 1)\simeq\mathbf 1$.
\end{lem}
We refer to the map $a_{\delta ,\nu}$ defined in the lemma as the \tdef{assembly map} associated with $\delta$ and $\nu$. Of particular interest is the case where $\cC = \pPrdual\base$, and \[A = \Sp_\base^{\ani{\tset}[\base]}, \quad B = \D{\tset},\quad \text{and} \quad C = \pcoShat{\base}{\tset}\] for some proper shape submersion $p\colon\tset\to\base$ between locally compact Hausdorff spaces. In this case, one can take $\nu$ to be the evaluation pairing \[\varepsilon\colon\pcoShat\base\tset\otimes\D\tset\too\Sp_\base,\] and $\delta$ to be the left adjoint of the composition
\[
  \D\tset\otimes\Sp_\base^{\ani{\tset}[\base]}\xtoo{\id\otimes\psi}\D\tset\otimes\D\tset\simeq\D{\tset\times\tset}\xtoo{\Delta^*}\D\tset\xtoo{X_*}\Sp_\base,
\]
where $\psi$ denotes the inclusion of locally constant sheaves into all sheaves. The associated assembly map
\[
  a_{\delta ,\varepsilon}\colon\pcoShat\base\tset\too\Sp_\base^{\ani{\tset}[\base]}
\]
is a categorification of the Waldhausen assembly map in the unparametrized case $\base=\pt$ by work of Bartels--Nikolaus, or more generally the parametrized Waldhausen assembly map of the Dwyer--Weiss--Williams index theorem by work of Volpe--Wolff--Ramzi. More precisely, for the parametrized K-theory functor \[\mathrm K\colon\pPrdual\base\too\Sp_\base,\] we have a natural identification of $\mathrm K(a_{\delta ,\varepsilon })$ with the parametrized Waldhausen assembly map as objects in the arrow category $\Fun (\Delta^1,\Sp_\base )$.

We will need a twisted variation on the above to deal with self-maps. In order to set this up, we need a basic property of the copairing described above:
\begin{lem}
  Let $p\colon \tset\to\base$ be a proper shape submersion between locally compact Hausdorff spaces, and let
  \[
    \delta_{\mathrm{BN}}\colon\Sp_\base\too\D\tset\otimes\Sp_\base^{\ani{\tset}[\base]}
  \]
  be the Bartels--Nikolaus copairing described in the discussion above. Given another proper shape submersion $q\colon Y\to\base$ and a map $f\colon \tset\to Y$ over $\base$, there is a natural identification
  \[
    (f^*\otimes\id )\delta_{\mathrm{BN}}\simeq (\id\otimes f_!)\delta_{\mathrm{BN}}
  \]
  of functors $\Sp_\base\to\D\tset\otimes\D Y$.
\end{lem}
\begin{proof}
  \Onote{This is shown in the proof of~\cite[Lem~3.4]{Harr2026b}, but the writing in that proof is horrendous, so we should write a clean version here.}.
\end{proof}
\begin{cons}\label{cons: twisted pairing and copairing}
  Let $(X,p\colon\tset\to\base ,f\colon\tset\to\tset )$ be a shapely dynamical system. Define a pairing \[\nu\colon (\pcoShat\base\tset ,f^?)\otimes (\D\tset ,f^*)\too (\Sp_\base ,\id)\] in $\Endtrl (\pPrl\base)$ to be the lax commutative square
  \[
    \begin{tikzcd}
      {\coShat_\base(\tset)\otimes\D\tset} & \Sp_\base\\
      {\coShat_\base(\tset)\otimes\D\tset} & \Sp_\base,
      \arrow["{\varepsilon}", from=1-1, to=1-2]
      \arrow["{f^?\otimes f^*}", from=1-1, to=2-1]
      \arrow["\id", from=1-2, to=2-2]
      \arrow[Rightarrow, from=2-1, to=1-2]
      \arrow["{\varepsilon}"', from=2-1, to=2-2]
    \end{tikzcd}
  \]
  where the 2-cell is the composite
  \[
    \varepsilon (f^?\otimes f^*)\simeq \varepsilon (\id\otimes f^*)(f^?\otimes \id)\simeq\varepsilon (f_?\otimes\id )(f^?\otimes\id)\simeq\varepsilon (f_?f^?\otimes\id)\too\varepsilon, 
  \]
  where the second natural equivalence is the definition of $f_?$.

  Similarly, we define a copairing
  \[
    \delta\colon (\Sp_\base ,\id)\too (\D\tset ,f^*)\otimes (\Sp_\base^{\ani{\tset}[\base]},f^*)
  \]
  to be the lax commutative square
  \[
    \begin{tikzcd}
      {\Sp_\base} & \D\tset\otimes\Sp_\base^{\ani{\tset}[\base]}\\
      {\Sp_\base} & \D\tset\otimes\Sp_\base^{\ani{\tset}[\base]},
      \arrow["{\delta_{\mathrm{BN}}}", from=1-1, to=1-2]
      \arrow["{\id}", from=1-1, to=2-1]
      \arrow["{f^*\otimes f^*}", from=1-2, to=2-2]
      \arrow[Rightarrow, from=2-1, to=1-2]
      \arrow["{\delta_{\mathrm{BN}}}"', from=2-1, to=2-2]
    \end{tikzcd},
  \]
  where $\delta_{\mathrm{BN}}$ is the Bartels--Nikolaus copairing described in the discussion preceding this construction, and the 2-cell is the composite
  \[
    (f^*\otimes f^*)\delta_{\mathrm{BN}}\simeq (\id\otimes f^*)(f^*\otimes\id)\delta_{\mathrm{BN}}\simeq (\id\otimes f^*)(\id\otimes f_!)\delta_{\mathrm{BN}}\simeq (\id\otimes f^*f_!)\delta_{\mathrm{BN}}\too \delta_{\mathrm{BN}},
  \]
  where the last map is the counit of the $f_!\dashv f^*$ adjunction.
\end{cons}
\begin{lem}\label{lem: twisted assembly is adjoint of assembly}
  Let
  \[
    \begin{tikzcd}
      \pcoShat\base\tset\arrow[d,"f_?"] \arrow[r,"a_{\delta ,\varepsilon}"] & \Sp_\base^{\ani{\tset}[\base]}\arrow[d,"f_!"]\\
      \pcoShat\base\tset\arrow[r,"a_{\delta ,\varepsilon }"] & \Sp_\base^{\ani{\tset}[\base]}
    \end{tikzcd},
  \]
  denote the commutative square coming from the naturality of $a_{\delta ,\varepsilon }$ as in~\cite[Thm~8.14]{Mukherjee_Nikolaus2026}. The lax commutative square defined by adjoining the vertical functors is canonically identified with the assembly map $a_{\delta ,\nu}$ associated with the pairing and copairing from~\cref{cons: twisted pairing and copairing}.
\end{lem}
\begin{proof}
  \Onote{unfold definitions}
\end{proof}
\begin{thm}
  Let $f\curvearrowright \tset\xrightarrow p\base$ be a shapely dynamical system. The lax commutative square
  \[
    \begin{tikzcd}
      \Sp_\base\arrow[d,equal] \arrow[r] & \Sp_\base^{\ani{\tset}[\base]}\arrow[d,"f^*"]\\
      \Sp_\base\arrow[r]\arrow[ru,Rightarrow, shorten <=5pt,shorten >=5pt] & \Sp_\base^{\ani{\tset}[\base]}
    \end{tikzcd}
  \]
  admits a factorization
  \begin{equation}
    \label{eq:categorical-factorization-of-param-reidemeister}
    \begin{tikzcd}
      \Sp_\base & {\coShat_\base(\tset)} & \Sp_\base^{\ani{\tset}[\base]} \\
      \Sp_\base & {\coShat_\base(\tset)} & \Sp_\base^{\ani{\tset}[\base]}
      \arrow["{w_{\tset}}", from=1-1, to=1-2]
      \arrow["\id"', from=1-1, to=2-1]
      \arrow["{a_{\delta ,\varepsilon} }", from=1-2, to=1-3]
      \arrow["{f^?}", from=1-2, to=2-2]
      \arrow["f^{\ast}", from=1-3, to=2-3]
      \arrow[Rightarrow, from=2-2, to=1-3]
      \arrow["{w_{\tset}}"', from=2-1, to=2-2]
      \arrow[Rightarrow, from=2-1, to=1-2]
      \arrow["{a_{\delta ,\varepsilon} }"', from=2-2, to=2-3]
    \end{tikzcd}
  \end{equation}
  natural in proper maps of shapely dynamical systems,
  where $w_\tset$ is the parametrized Wall object and the rightmost square is the assembly map in $\Endtrl (\pPrlf{B})$ associated with the $\delta$ and $\nu$ constructed in~\cref{cons: twisted pairing and copairing}. 
\end{thm}
\begin{proof}
  The left-hand 2-cell is adjoint to the natural map \[\theta\colon f_?w_\tset\to w_\tset\]  described in Construction 3.1 of~\cite{Harr2026b}. We can therefore describe this 2-cell as
  \[
    w_\tset\xtoo{\text{unit}} f^?f_?w_\tset\xtoo{f^?(\theta )} f^?w_\tset.
  \]
  By~\cref{lem: twisted assembly is adjoint of assembly}, the right-hand 2-cell is the Beck--Chevalley morphism
  \[
    a_{\delta ,\varepsilon }f^?\xtoo{\text{unit}} f^*f_!a_{\delta ,\varepsilon }f^?\simeq f^*a_{\delta ,\varepsilon }f_?f^?\xtoo{\text{counit}}f^*a_{\delta ,\varepsilon },
  \]
  where the equivalence is the naturality 2-cell of~\cite[Thm~8.14]{Mukherjee_Nikolaus2026}. By functoriality of unit transformations, we find that the composed 2-cell is adjoint to the unit map \[f_!f^*a_{\delta ,\varepsilon }w_\tset\too a_{\delta ,\varepsilon },\] and~\cite[Lem~3.4]{Harr2026b} gives a natural identification of this map with the counit map
  \[
    f_!\SS_\tset\too\SS_\tset.
  \]
  But this is exactly the definition of the 2-cell~\labelcref{eq:reidemeister lax square}, see~\cref{subsec: homotopical lefschetz}. 
\end{proof}
\begin{prop}\label{prop: lax calculation of trace on coshv hat}
  Let $f\curvearrowright \tset\xrightarrow p\base$ be a shapely dynamical system. 
  The map
  \[
    \TrcL (a_\delta)\colon \TrcL (f^?\curvearrowright\coShat_\base (\tset ))\too \TrcL(f^*\curvearrowright \Sp_\base^{\ani{\tset}[\base]})\simeq \SS [\cL^{f}\ani{\tset}[\base]]
  \]
  admits a factorization
  \begin{equation}
    \label{eq:lax-trace-on-cosheaves}
    \begin{tikzcd}
      &  \pbmsec{X^f}{\base}\\
      \TrcL (f^?\curvearrowright\coShat_\base (\tset )) & {\SS [\cL^{f}\ani{\tset}[\base]]},
      \arrow[from=1-2, to=2-2]
      \arrow[from=2-1, to=1-2]
      \arrow[from=2-1, to=2-2]
    \end{tikzcd}
  \end{equation}
  in $\Sp_\base$, natural in proper maps of shapely dynamical systems. Furthermore, the composite
  \[
    \SS_\base\simeq\TrcL (\id\curvearrowright\Sp_\base)\to \TrcL (f^?\curvearrowright\coShat_\base (\tset ))\to \pbmsec{X^f}{\base}
  \]
  is homotopic to $\lambda_f$.
\end{prop}
\begin{proof}
  Aside from the last statement, everything follows immediately from~\cref{lem: bartels-nikolaus nonsense}. The last statement on the other hand is the identification of the parametrized THH-assembly with the result of applying THH to the assembly functor. \Onote{details or reference for the last bit}
\end{proof}

\subsubsection{The Nielsen fixed point theorem}\label{subsec:classical-nielsen}
We now specialize the results of the previous section to the case
where $\base = \pt$. 
Let $\tset$ be a topological set that is compact Hausdorff and locally of constant shape. For a self-map $f\colon \tset\to \tset$, the so-called Nielsen fixed point theorem of Reidemeister~\cite{Reidemeister1936} and Wecken~\cite{Wecken1941a,Wecken1941b,Wecken1942} gives a lower bound on the number of fixed points of $f$ that only depends on the homotopy class of $f$. The lower bound is expressed in terms of the following numerical invariant:
\begin{defn}\label{defn: Nielsen number}
  The \tdef{Nielsen number} of $f$, denoted $\mdef{N_f}$, is the number of nonzero coordinates of the Reidemeister character $\rho_f\in\pi_0\SS\lbrack\cL^f\ani{\tset}\rbrack$ with respect to the standard basis
  \begin{equation}
    \label{eq:standard-basis-of-reidemeister-target}
    \pi_0\SS\lbrack\cL^f\ani{\tset}\rbrack\simeq\ZZ\lbrack\pi_0\cL^f\ani{\tset}\rbrack.
  \end{equation}
\end{defn}
By construction, $N_f$ depends only on the homotopy class of $f$.\footnote{In other approaches to Nielsen fixed point theory, the Nielsen number---or rather the Reidemeister trace---is often defined in terms of $\tset^f$, and one then has to do some work to show that it is homotopy-invariant. For us, the homotopy-invariance of the Nielsen number is congenital, and the work that we do instead goes towards showing that it relates to $\tset^f$, which is the content of~\cref{thm: lefschetz-nielsen}.}
\begin{prop}[Nielsen fixed point theorem]
  $\#\tset^f \geq N_f\,$.
\end{prop}
Since the Lefschetz trace is the sum of the coordinates of the Reidemeister trace, we have
\begin{equation}
  \label{eq:lefschetz-vs-nielsen-fixed-pt}
  \text{Lefschetz trace }\geq 1\Rightarrow N_f\geq 1.
\end{equation}
We can therefore view the Nielsen fixed point theorem as a refinement of the Lefschetz fixed point theorem. Note however that aside from~\eqref{eq:lefschetz-vs-nielsen-fixed-pt}, there is no general formula relating the Nielsen number and the Lefschetz trace.
\begin{proof}
  We may assume that $\# \tset^f < \infty$, as the statement is trivial otherwise. Consider the commutative diagram
  \[
    \begin{tikzcd}
      \tset^f\arrow[d,"\id"]\arrow[r,"i"] & \tset\arrow[d,"f"] \\
      \tset^f\arrow[r,"i"] & \tset,
    \end{tikzcd}
  \]
  where $i$ denotes the inclusion. By~\cref{thm: lefschetz-nielsen}, we get a commutative diagram
  \[
    \begin{tikzcd}
      \SS\lbrack \tset^f\rbrack \arrow[d] \arrow[r,"\sim"] & \pbmsec{X^f}{\base}\arrow[d,"\mathfrak a"] \\
      \SS\lbrack\cL \tset^f\rbrack\arrow[r] & \SS\lbrack\cL^f\tset\rbrack,
    \end{tikzcd}
  \]
  where the left vertical map is induced by the canonical map $\tset^f\to\cL \tset^f$, which is an equivalence since $\tset^f$ is discrete. On $\pi_0$, we thus find that the image of $\mathfrak a\colon \pbmsec{X^f}{\base}\to \SS\lbrack\cL^f\tset\rbrack$ is contained in the span of the basis vectors given by path components of $(x,c_x)\in\cL^f\tset$, where $x\in \tset^f$ and $c_x$ is the constant loop at $x$. It follows that $N_f\leq\# \tset^f$ as desired.
\end{proof}

% \subsection{Another trace ``calculation''}
% \todo{Should this still be hear?}
% Keeping with the pattern laid out in~\cref{sec: lefschetz theory}, the next step in our treatment of Lefschetz--Nielsen theory is a calculation of the categorical trace that shows up as the apex in the triangle~\eqref{eq:nielsen-triangle-without-identification}. Unlike in the Lefschetz and Reidemeister trace incarnations of this step, we will not be able to calculate $\tr{f^?\curvearrowright\coShat (\tset )}{}$ in the strict sense (without a big effort); instead, we will give a lax calculation that suffices for the purposes of fixed point theory. This idea is again adapted to our setting from unpublished work of Bartels and Nikolaus, who show using simple manipulations in closed monoidal categories that the map induced on $K$-theory by their assembly functor $a_\delta\colon\coShat (\tset )\to \Sp^{\fundgroupoid (\tset )}$ factors through $\tset_*\tset^!K\SS$, thereby removing the difficult~\cref{thm: K theory of completed cosheaves} from the dependencies of the customer-facing~\cref{thm: Bartels--Nikolaus geometric topology}.

% We should expect the calculation of $\tr{f^?\curvearrowright\coShat (\tset )}{}$ to be difficult, because for $f=\id$ it is already a difficult theorem of Efimov.\todo{motivate linearizing over free algebra on dualizable generator}
% \begin{defn}[The free algebra on one dualizable object]
%   $\freeondual = \Fun (\mathrm{Bord}^1_{\mathrm{fr}}(\RR^1)^{\op},\Sp )$\todo{define Hopf algebra structure}
% \end{defn}

%%% Local Variables:
%%% mode: LaTeX
%%% TeX-master: "main"
%%% End:
